\documentclass[tikz,border=2mm]{standalone}

\usepackage{amsmath}
\usetikzlibrary{arrows,positioning}
\usetikzlibrary{fit,backgrounds}
\usepackage{xcolor}

\begin{document}

\definecolor{beige}{RGB}{245,245,220}
\colorlet{pdcol}{red!12}
\colorlet{hdcol}{green!14}
\colorlet{shcol}{blue!10}
\colorlet{hmcol}{orange!12}

\begin{tikzpicture}[node distance=0.5cm,auto,>=latex',every node/.append style={align=center},int/.style={draw, shape=circle, minimum size=1.2cm},inverter/.style={circle,draw,inner sep=2pt,minimum size=6mm},posnode/.style={shape=circle, draw=orange, line width=2},
  negnode/.style={shape=circle, draw=black, line width=2}]
  
  \node(eq1)[draw, shape=rectangle,inner sep=10pt, line width=2,  rounded corners=6pt,
  text width=0.4\linewidth,
  fill=beige,
  align=center,
  minimum height=2.4cm] at (0,0) {
$\displaystyle
  \frac{dx_i}{dt} = x_i(\pi_i(\mathbf{x}) - \bar{\pi}(\mathbf{x}))
  $
};
  \node (title1) at (0,1.6) {\textbf{1. Replicator dynamics}};

  \node(eq2)[draw, shape=rectangle,inner sep=10pt, line width=2,  rounded corners=6pt,
  text width=0.4\linewidth,
  fill=beige,
  align=center,
  minimum height=2.4cm] at (7,0) {
  At $\mathbf{e}_i$, can $j$ invade?\\
  \vspace{0.2em}
$\displaystyle
  \pi_j(\mathbf{e}_i) > \pi_i(\mathbf{e}_i)
  $\\
  \vspace{0.2em}
  At $\mathbf{e}_j$, can $i$ invade?\\
  \vspace{0.2em}
  $\displaystyle
  \pi_i(\mathbf{e}_j) > \pi_j(\mathbf{e}_j)
  $
};
  \node (title2) at (7,1.6) {\textbf{2. Mutual invasion conditions}};
  
  \node(eq3)[draw, shape=rectangle,inner sep=10pt, line width=2,  rounded corners=6pt,
  text width=0.4\linewidth,
  fill=beige,
  align=center,
  minimum height=2.4cm] at (14,0) {
  $\displaystyle
  W_i(\mathbf{x}) = \exp(\pi_i(\mathbf{x}))
  $\\
  \vspace{0.8em}
  Preserves invasion conditions and allows ratios
};
  \node (title3) at (14,1.6) {\textbf{3. Multiplicative payoffs}};
  
  
  \node(eq4)[draw, shape=rectangle,inner sep=10pt, line width=2,  rounded corners=6pt,
  text width=0.4\linewidth,
  fill=beige,
  align=center,
  minimum height=3.5cm] at (0,0-4.5) {
$\displaystyle
\begin{aligned}
  1- \rho &= 1-\sqrt{\frac{W_i(\mathbf{e}_i) W_j(\mathbf{e}_j)}{W_j(\mathbf{e}_i) W_i(\mathbf{e}_j)}}\\
  \frac{\kappa_j}{\kappa_i} &= \sqrt{\frac{W_j(\mathbf{e}_i)W_j(\mathbf{e}_j)}{W_i(\mathbf{e}_i) W_i(\mathbf{e}_j)}}
  \end{aligned}
  $
};
  \node (title4) at (0,2-4) {\textbf{4. SCT decomposition}};


  \node(eq5)[draw, shape=rectangle,inner sep=10pt, line width=2,  rounded corners=6pt,
  text width=0.4\linewidth,
  fill=beige,
  align=center,
  minimum height=3.5cm] at (7,0-4.5) {
$\displaystyle
\rho < \frac{\kappa_j}{\kappa_i} < \frac{1}{\rho}
  $\\
  \vspace{0.8em}
  Stabilizing niche differences must overcome average fitness differences
};
  \node (title5) at (7,2-4) {\textbf{5. Coexistence criterion}};


\begin{scope}[shift={(14,-4.5)}, x=3cm, y=2.52cm]
  \def\xmin{-0.6}
  \def\xmax{ 0.6}
  \def\ymin{-0.7}
  \def\ymax{ 0.7}

  \pgfmathsetmacro{\xcrit}{1 - exp(-\ymax)}


            \fill[pdcol]
            (\xmin,\ymax) --
            (\xmax,\ymax) --
            plot[domain=\xmax:0, samples=150]
                (\x,{min(-ln(1-\x),\ymax)}) --
            plot[domain=0:\xmin, samples=150]
                (\x,{min( ln(1-\x),\ymax)}) --
            cycle;

            \fill[hmcol]
            (\xmin,\ymin) --
            (\xmax,\ymin) --
            plot[domain=\xmax:0, samples=150]
                (\x,{max( ln(1-\x),\ymin)}) --
            plot[domain=0:\xmin, samples=150]
                (\x,{max(-ln(1-\x),\ymin)}) --
            cycle;
            
             \fill[shcol]
            plot[domain=\xmin:0, samples=150]
                (\x,{ ln(1-\x)}) --
            plot[domain=0:\xmin, samples=150]
                (\x,{-ln(1-\x)}) --
            cycle;
            
            \fill[hdcol]
            plot[domain=0:\xmax, samples=150]
                (\x,{max( ln(1-\x),\ymin)}) --
            plot[domain=\xmax:0, samples=150]
                (\x,{min(-ln(1-\x),\ymax)}) --
            cycle;


  % region labels
  \node[font=\footnotesize, align=center, text width=1.7cm] at (-0.04,0.43)
    {$j$ excludes $i$};

  \node[font=\footnotesize, align=center, text width=1.7cm] at (0.35,0.00)
    {Coexist-\\ence};

  \node[font=\footnotesize, align=center, text width=1.7cm] at (-0.35,0.00)
    {Priority\\effect};

  \node[font=\footnotesize, align=center, text width=1.7cm] at (-0.04,-0.43)
    {$i$ excludes $j$};
    
  % border
  \draw[line width=0.8pt] (\xmin,\ymin) rectangle (\xmax,\ymax);

  % axes
  \draw[->, line width=0.8pt] (\xmin,\ymin) -- (\xmax+0.09,\ymin);
  \draw[->, line width=0.8pt] (\xmin,\ymin) -- (\xmin,\ymax+0.12);

  % axis labels
  \node[below=10pt] at (0,-0.6)
    {\footnotesize Niche difference, $1-\rho$};

  \node[rotate=90, align=center] at (\xmin-0.16,0)
    {\footnotesize Fitness difference, $\log(\kappa_j/\kappa_i)$};

  % boundary curves
  \draw[line width=1pt]
    plot[domain=\xmin:\xcrit, samples=150]
      (\x,{ln(1-\x)});

  \draw[line width=1pt]
    plot[domain=\xmin:\xcrit, samples=150]
      (\x,{-ln(1-\x)});
\end{scope}

\node (title6) at (14,2-4) {\textbf{6. Niche-fitness difference space}};

\end{tikzpicture}
\end{document}